% !TEX root = ../main.tex

\chapter{Diseño de Solución}
\label{chap:diseno-solucion}

A partir del problema formulado en la Sección~\ref{sec:problema-formulacion}, se observa que, aunque \texttt{ApplicativeDo} puede mejorar el plan de ejecución completamente secuencial, la preservación del orden relativo de los statements limita el espacio de planes que puede explorar. Bajo una hipótesis explícita de conmutatividad, aplicar \texttt{ApplicativeDo} sobre un reordenamiento válido de statements puede exponer un plan de ejecución de menor costo que el obtenido a partir del orden original \cite[pp.~94--95]{MarlowPJKM2016}.

La solución propuesta por la memoria consiste en extender el espacio de búsqueda de \texttt{ApplicativeDo}: para un bloque \texttt{do} declarado como conmutativo, se generan los reordenamientos válidos, se aplica el algoritmo existente sobre cada uno y se selecciona el plan de ejecución de menor costo. De este modo, la extensión aprovecha la libertad adicional del bloque \texttt{do} sin reemplazar \texttt{ApplicativeDo} ni modificar el tratamiento de los bloques ordinarios.

En este capítulo se desarrolla el diseño de dicha solución. Primero, se presenta el flujo general de la extensión, desde la declaración del bloque como conmutativo hasta la selección del plan de ejecución de menor costo. A continuación, se introduce el \textit{grafo de precedencia RAW} como representación de las restricciones entre los statements de un bloque \texttt{do}, junto con el modelo de dependencias que justifica por qué basta considerar dependencias locales de tipo \textit{Read-after-Write}. Finalmente, se caracterizan los reordenamientos válidos como ordenamientos topológicos del grafo y se describe el recorrido exhaustivo utilizado para generarlos.

% !TEX root = ../../main.tex

\section{Flujo general de la extensión}
\label{sec:solucion-flujo}

El diseño propuesto amplía el flujo habitual de \texttt{ApplicativeDo} al incorporar la exploración de órdenes alternativos antes de construir el plan de ejecución. Sobre un bloque \texttt{do} habilitado, la solución representa las dependencias locales mediante un grafo de precedencia RAW. Sus aristas determinan qué relaciones deben conservarse y, por tanto, caracterizan como reordenamientos válidos a las secuencias completas que respetan todas esas precedencias. Dentro del alcance experimental de esta memoria, la hipótesis explícita de conmutatividad permite considerar dichos reordenamientos semánticamente equivalentes al orden original; la ausencia de una dependencia local, por sí sola, no sería suficiente para autorizar el intercambio sobre una mónada arbitraria \cite[p.~94]{MarlowPJKM2016}.

Cada reordenamiento válido se entrega a la etapa de \textit{rearrange} de \texttt{ApplicativeDo}. Como se explicó en la Sección~\ref{sec:applicative-do-rearrange}, esta etapa conserva el orden de statements recibido y construye un árbol que combina secuencias monádicas y agrupaciones aplicativas. La extensión compara los árboles obtenidos para todos los reordenamientos mediante el modelo de costo de la Expresión~\ref{expr:modelo-costos-estructural} y selecciona el primer candidato de menor costo. El árbol seleccionado continúa luego hacia la etapa de \textit{desugaring}. Por consiguiente, la contribución amplía el espacio de búsqueda, mientras que \texttt{ApplicativeDo} sigue siendo responsable de construir el plan asociado a cada reordenamiento. La Figura~\ref{fig:solucion-pipeline-conceptual} resume este proceso.

\begin{figure}[ht]
\centering
\vspace{0.3cm}
\captionsetup{skip=15pt}
\begin{tikzpicture}[
    node distance=1.05cm and 0.45cm,
    solution step/.style={
        draw=codeframe,
        fill=codebg,
        rounded corners=2pt,
        minimum width=3.25cm,
        minimum height=1.05cm,
        text width=2.9cm,
        align=center,
        font=\small
    },
    solution contribution/.style={
        solution step,
        draw=codekeyword,
        fill=codekeyword!18!codebg
    },
    solution arrow/.style={precedence edge,draw=codekeyword,line width=0.7pt}
]
    \node[solution step] (marked) {Bloque \texttt{do} declarado como conmutativo};
    \node[solution contribution,right=of marked] (orders) {Generación de reordenamientos válidos};
    \node[solution step,right=of orders] (ado) {Aplicación de \texttt{ApplicativeDo} sobre cada reordenamiento};
    \node[solution step,right=of ado] (selected) {Selección del plan de ejecución de menor costo};

    \draw[solution arrow] (marked) -- (orders);
    \draw[solution arrow] (orders) -- (ado);
    \draw[solution arrow] (ado) -- (selected);
\end{tikzpicture}
\caption{Flujo general de la extensión propuesta.}
\label{fig:solucion-pipeline-conceptual}
\end{figure}

El diseño separa así dos responsabilidades. El programador afirma que el orden de efectos independientes puede intercambiarse en el bloque \texttt{do} marcado. El compilador restringe la exploración a reordenamientos que conservan las relaciones locales de definición y uso, evalúa los planes de ejecución correspondientes y selecciona el menos costoso, sin modificar la compilación de los bloques que no son declarados como conmutativos.

En este flujo, el núcleo algorítmico de la extensión corresponde a la generación de reordenamientos válidos, cuyo diseño se divide en dos fases. La primera construye el grafo de precedencia RAW que representa las restricciones entre los statements del bloque; la segunda recorre exhaustivamente dicho grafo para generar los reordenamientos que preservan esas restricciones. Ambas fases se desarrollan, respectivamente, en las Secciones~\ref{sec:solucion-grafo-precedencia-raw} y~\ref{sec:solucion-generacion-reordenamientos}.

% !TEX root = ../../main.tex

\section{Grafo de precedencia RAW}
\label{sec:solucion-grafo-precedencia-raw}

En la primera fase del proceso de generación de reordenamientos válidos se requiere una representación capaz de expresar las precedencias que deben conservarse entre los statements de un bloque \texttt{do} habilitado. Siguiendo el uso de grafos de dependencias para formular e implementar optimizaciones de compiladores \cite{KuckEtAl1981}, la solución utiliza un grafo de precedencia cuyas aristas representan exclusivamente dependencias locales de tipo \textit{Read-after-Write}; este grafo especializado se denomina \textit{grafo de precedencia RAW}. La Figura~\ref{fig:main-example-precedence-graph} presentó previamente esta representación para el programa del Código~\ref{lst:main-example}: cada vértice corresponde a un statement y cada arista conecta un statement productor con otro que consume uno de sus valores. Sin embargo, antes de definir formalmente el grafo, es necesario justificar por qué basta considerar dependencias RAW y por qué las dependencias \textit{Write-after-Read} y \textit{Write-after-Write} no originan precedencias adicionales entre statements renombrados. A continuación, se presenta el modelo de dependencias que fundamenta esta decisión.

\subsection{Modelo de dependencias}
\label{sec:solucion-modelo-dependencias}

El punto de partida es un modelo general de dependencias entre lecturas y escrituras. Considérense dos statements $s_i$ y $s_j$ tales que $i<j$. Cada statement $s_k$ posee un conjunto $R_k$ de variables leídas y un conjunto $W_k$ de variables escritas. A partir de estos conjuntos se distinguen tres clases de dependencia: \textit{Read-after-Write} (RAW), \textit{Write-after-Read} (WAR) y \textit{Write-after-Write} (WAW):

\begin{ReportExpression}{Dependencias de escritura y lectura de variables entre dos statements $s_i$ y $s_j$.}{expr:dependencias-raw-war-waw}
\[
\begin{aligned}
\operatorname{RAW}(i,j) &\iff W_i\cap R_j\neq\emptyset,\\
\operatorname{WAR}(i,j) &\iff R_i\cap W_j\neq\emptyset,\\
\operatorname{WAW}(i,j) &\iff W_i\cap W_j\neq\emptyset.
\end{aligned}
\]
\end{ReportExpression}

Una dependencia RAW establece una relación productor--consumidor: $s_j$ lee una variable cuyo valor fue escrito por $s_i$. Una dependencia WAR aparece cuando $s_i$ debe leer una variable antes de que $s_j$ la sobrescriba; invertir ambos statements haría que la lectura observara el nuevo valor. Finalmente, una dependencia WAW relaciona dos escrituras sobre una misma variable y obliga a conservar cuál de ellas determina su valor posterior. Estas dos últimas relaciones presuponen que distintas escrituras pueden actuar sobre una misma identidad mutable, como ocurre en un modelo imperativo de asignaciones. Alterar cualquiera de estas precedencias puede modificar los valores observados o el estado final del programa \cite{Bernstein1966,ErtlKrall1996}.

La \textit{do-notation} de \textsf{Haskell} presenta un comportamiento diferente. En un statement de la forma \verb|p <- e|, el patrón \verb|p| introduce una nueva vinculación léxica, denominada \textit{binder}; no actualiza destructivamente una variable definida con anterioridad. Si un statement posterior reutiliza el mismo identificador textual, introduce otro binder que oculta al primero en su nuevo alcance. Este comportamiento concuerda con la pureza descrita en la Sección~\ref{sec:fundamentos-haskell} y con las reglas de alcance de la \textit{do-notation} \cite[sección~3.14]{Haskell2010}.

El Renamer de \textsf{GHC} hace explícita esta diferencia al resolver los identificadores del programa. Cada binder recibe un \texttt{Name} único y cada uso queda asociado con la identidad que se encuentra en alcance en ese punto. Por ello, dos definiciones escritas como \texttt{x} en el programa fuente no representan dos escrituras internas sobre el mismo nombre \cite{Appel1998,GHCRenamerSources2026}.

El Código~\ref{lst:original-vs-renaming} ilustra este comportamiento con un bloque \texttt{do} de tipo \texttt{Maybe Int}. En el programa fuente, \texttt{y} lee la primera vinculación de \texttt{x} antes de que un statement posterior vuelva a utilizar ese identificador. La columna derecha presenta esquemáticamente las identidades diferenciadas durante el renombrado del compilador:

\noindent\begin{minipage}{\textwidth}
\begin{lstlisting}[style=haskell,caption={Representación esquemática de los binders antes y después del renombrado que ejecuta \textsf{GHC} a todo programa.},label={lst:original-vs-renaming}]
    do                                      do
      x <- Just 2                             x_a1T2 <- Just 2
      y <- Just ((2*) x)                      y_a1T3 <- Just ((2*) x_a1T2)
      x <- Just 3              ==>            x_a1T4 <- Just 3
      return (x + y)                          return (x_a1T4 + y_a1T3)
\end{lstlisting}
\end{minipage}

La primera y la tercera definición escriben, respectivamente, los nombres \texttt{x\_a1T2} y \texttt{x\_a1T4}. Como son identidades diferentes, sus conjuntos de escritura no se intersectan y, por tanto, no existe una dependencia WAW. De manera análoga, el segundo statement lee \texttt{x\_a1T2}, mientras que el tercero escribe \texttt{x\_a1T4}; la intersección entre esa lectura y la escritura posterior también es vacía, por lo que no existe una dependencia WAR. La relación local efectiva entre los statements reordenables es RAW: la definición de \texttt{x\_a1T2} debe permanecer antes del statement que define \texttt{y\_a1T3}. En consecuencia, el grafo de precedencia utilizado por la solución se construye exclusivamente a partir de dependencias RAW entre nombres renombrados.

Sea $S=(s_1,\ldots,s_n)$ la secuencia de statements que participan en el reordenamiento. El statement terminal, representado internamente como \texttt{LastStmt}, no pertenece a $S$ porque produce el resultado del bloque y permanece fijo \cite[p.~95]{MarlowPJKM2016}. La anotación de variables libres que entrega el Renamer para el cuerpo de cada statement contiene los \texttt{Name}s utilizados que no son introducidos por ese mismo cuerpo \cite{GHCRenamerSources2026}. Esto incluye tanto referencias a binders definidos por statements anteriores como nombres globales, entre ellos funciones, operadores y constructores de datos. Por ejemplo, en el segundo statement del Código~\ref{lst:original-vs-renaming} aparecen como nombres libres \texttt{x\_a1T2}, el constructor \texttt{Just} y el operador \verb|(*)|, pero solo el primero corresponde a un valor producido dentro del bloque \texttt{do}. Para distinguir las lecturas locales de las referencias externas se definen los siguientes conjuntos para cada statement $s_i$:

\begin{ReportExpression}{Escrituras y lecturas locales de un statement renombrado.}{expr:lecturas-escrituras-locales}
\[
\begin{aligned}
\operatorname{WRITE}(s_i) &= \operatorname{binders}(s_i),\\
\operatorname{ALL\_WRITES} &= \bigcup_{k=1}^{n}\operatorname{WRITE}(s_k),\\
\operatorname{READ}_{\mathrm{local}}(s_i)
  &= \operatorname{FV}(s_i)\cap\operatorname{ALL\_WRITES}.
\end{aligned}
\]
\end{ReportExpression}

El conjunto $\operatorname{WRITE}(s_i)$ contiene los \texttt{Name}s de los binders introducidos por $s_i$, mientras que $\operatorname{ALL\_WRITES}$ reúne todos los nombres producidos por los statements reordenables del bloque \texttt{do}. Por su parte, $\operatorname{FV}(s_i)$ conserva el conjunto completo de nombres libres informado por el Renamer de \textsf{GHC}. Al intersectarlo con $\operatorname{ALL\_WRITES}$, $\operatorname{READ}_{\mathrm{local}}(s_i)$ retiene únicamente los nombres cuyo valor se origina por un binder de un statement anterior. En el ejemplo del Código~\ref{lst:original-vs-renaming}, \texttt{x\_a1T2} permanece como lectura local del segundo statement, mientras que \texttt{Just} y \verb|(*)| se descartan de $\operatorname{READ}_{\mathrm{local}}(s_2)$ porque, aunque pertenecen a $\operatorname{FV}(s_2)$, no forman parte de $\operatorname{ALL\_WRITES}$.

Con esta información, la restricción relevante entre dos statements $s_i$ y $s_j$ queda dada por

\begin{ReportExpression}{Dependencia RAW local utilizada por la extensión.}{expr:dependencia-raw-local}
\[
i\longrightarrow j
\quad\iff\quad
i<j
\quad\land\quad
\operatorname{WRITE}(s_i)
\cap
\operatorname{READ}_{\mathrm{local}}(s_j)
\neq\emptyset.
\]
\end{ReportExpression}

La condición $i<j$ orienta la dependencia según el orden del programa original. La intersección no vacía indica que existe al menos un \texttt{Name} definido por $s_i$ y utilizado localmente por $s_j$; dicho nombre actúa como testigo de una relación productor--consumidor que debe preservarse \cite{HarroldSoffa1994}. Debido a que la comparación se realiza sobre nombres renombrados, la arista no puede surgir solo por la reutilización de un mismo identificador textual.

El grafo de precedencia RAW de un bloque \texttt{do} se define entonces como:

\begin{ReportExpression}{Grafo de precedencia RAW inducido por las dependencias locales de un bloque \texttt{do}.}{expr:grafo-precedencia-raw}
\[
G_{\mathrm{RAW}}=(V,E_{\mathrm{RAW}}),
\qquad
V=\{1,\ldots,n\},
\qquad
E_{\mathrm{RAW}}=\{(i,j)\in V\times V\mid i\longrightarrow j\}.
\]
\end{ReportExpression}

Cada vértice representa un statement reordenable y cada arista registra una dependencia RAW local que, al mismo tiempo, impone una precedencia productor--consumidor. Como toda arista satisface $i<j$, los índices aumentan estrictamente a lo largo de cualquier camino y $G_{\mathrm{RAW}}$ es acíclico. El conjunto $E_{\mathrm{RAW}}$ contiene todas las dependencias detectadas por el criterio anterior y no corresponde necesariamente a una reducción transitiva; la relación de alcanzabilidad del grafo induce el orden parcial que debe respetar todo reordenamiento válido \cite{Pratt1986}.

El grafo de precedencia RAW expresa así las restricciones estructurales que determinan la validez de un reordenamiento; por sí solo no garantiza equivalencia semántica, que en esta memoria depende de la hipótesis explícita de conmutatividad \cite[secciones~5 y~9]{Kock2012}. A partir del orden parcial inducido por este grafo, la segunda fase determina los reordenamientos de statements que preservan todas las precedencias.

% !TEX root = ../../main.tex

\section{Generación de reordenamientos válidos}
\label{sec:solucion-generacion-reordenamientos}

Una vez construido el grafo de precedencia RAW $G_{\mathrm{RAW}}$ asociado a un bloque \texttt{do}, generar los reordenamientos válidos equivale a obtener los ordenamientos topológicos del grafo. Un ordenamiento topológico es una forma de disponer todos sus vértices en una secuencia que respeta la dirección de las aristas \cite{Kahn1962}. En este caso, si existe una arista $i\to j$, el statement $s_i$ debe aparecer antes que $s_j$; si el grafo no impone una precedencia entre dos statements, estos pueden aparecer en cualquier orden relativo. Esta caracterización es central para el diseño de la solución: transforma la preservación de las dependencias locales en una propiedad estructural verificable sobre el grafo. En otras palabras, para determinar si un reordenamiento es válido basta comprobar que el origen de cada arista aparezca antes que su destino. Como toda arista satisface $i<j$, el orden original $(s_1,\ldots,s_n)$ también es topológico y permanece entre los candidatos efectivamente evaluados.

El procedimiento construye estos ordenamientos una posición a la vez. En cada paso puede seleccionar cualquier statement que todavía no haya sido agregado y cuyos predecesores ya aparezcan en la secuencia. Si existen varios statements disponibles, el procedimiento explora cada elección como una alternativa distinta. Este proceso se repite hasta incorporar todos los statements y, al considerar todas las elecciones posibles, se obtienen todos los ordenamientos topológicos del grafo y por ende todos los reordenamientos válidos del bloque \texttt{do} \cite{KnuthSzwarcfiter1974}.

Aplicado al programa del Código~\ref{lst:main-example}, este procedimiento debe preservar las tres aristas del grafo de precedencia RAW presentado en la Figura~\ref{fig:main-example-precedence-graph}. Los cinco ordenamientos topológicos que satisfacen estas restricciones se presentan en la Tabla~\ref{tab:solucion-candidatos-ejemplo}.

\begin{table}[ht]
\centering
\TableSetup
\TableFrame{%
\begin{tabular}{cc}
\rowcolor{tableHeaderBg}
\TableHead{Índice} & \TableHead{Ordenamiento topológico} \\
0 & \texttt{[1,2,3,4]} \\
1 & \texttt{[1,3,2,4]} \\
2 & \texttt{[1,3,4,2]} \\
3 & \texttt{[3,1,2,4]} \\
4 & \texttt{[3,1,4,2]} \\
\end{tabular}%
}
\caption{Ordenamientos topológicos del grafo de precedencia RAW de la Figura~\ref{fig:main-example-precedence-graph}.}
\label{tab:solucion-candidatos-ejemplo}
\end{table}

El primer ordenamiento conserva la secuencia original; los demás intercalan statements independientes sin invertir ninguna de las tres aristas del grafo.

Así, el grafo de precedencia RAW determina las restricciones de un bloque \texttt{do} y su recorrido exhaustivo materializa los reordenamientos válidos que se entregan a la etapa de \textit{rearrange} de \texttt{ApplicativeDo}. El Capítulo~\ref{chap:implementacion-solucion} desarrolla la implementación de ambas fases dentro del Renamer de \textsf{GHC}.

